\documentclass[tikz,border=20pt]{standalone}
\usepackage{ctex}
\usepackage{tikz}
\usetikzlibrary{arrows.meta, positioning, fit, backgrounds, calc}

\begin{document}
\begin{figure}[htbp]
	\centering % 核心：确保水平居中
	% 核心：resizebox 确保图像宽度正好等于正文宽度 (\textwidth)
	% 如果觉得太满，可以改为 0.9\textwidth
	\resizebox{\textwidth}{!}{ 
		\begin{tikzpicture}[
			font=\small\sffamily,
			>=Stealth,
			box/.style={
				rectangle, draw=black!80, fill=white, very thick, 
				minimum width=42mm, minimum height=12mm, rounded corners=2pt, align=center
			},
			highlight/.style={box, draw=blue!70!black, fill=blue!5},
			route/.style={->, thick, draw=black!70},
			sync/.style={->, dashed, thick, draw=black!40}
			]
			
			% --- 1. 核心布局 (垂直中轴严格对齐) ---
			\node[box] (gptp) at (0, 4.5) {gPTP 时间同步\\Global Time Base};
			\node[highlight] (cpu) at (0, 0) {中央处理器\\Central Processor};
			\node[box] (ring) at (0, -4) {环形存储\\Ring Buffer};
			
			% 左右对称
			\node[box] (can) at (-5.8, 1.5) {CAN 总线\\CAN Frame Stream};
			\node[box] (eth) at (-5.8, -1.5) {以太网监控\\Ethernet Traffic Monitor};
			\node[box] (ids) at (5.8, 0) {CAN ClockIDS\\实时检测引擎};
			
			% --- 2. 背景装饰 ---
			\begin{scope}[on background layer]
				\node[fill=gray!5, rounded corners=8pt, fit=(gptp)(can)(eth)(cpu)(ids)(ring), inner sep=10mm] (bg) {};
			\end{scope}
			
			% --- 3. 信号连线 (简约化处理) ---
			
			% 输入流：折线接入，避免交叉
			\draw[route] (can.east) -| ([xshift=-8mm, yshift=2mm]cpu.west) -- ([yshift=2mm]cpu.west)
			node[pos=0.2, above, font=\scriptsize] {CAN 数据};
			\draw[route] (eth.east) -| ([xshift=-8mm, yshift=-2mm]cpu.west) -- ([yshift=-2mm]cpu.west)
			node[pos=0.2, below, font=\scriptsize] {以太网上下文};
			
			% 检测流
			\draw[route] (cpu.east) -- (ids.west)
			node[midway, above, font=\scriptsize] {特征提取};
			
			% 存储双向流：垂直并行，直接向上/下，绝不拐弯
			% 写入 (偏移 +8mm)
			\draw[route] ([xshift=8mm]cpu.south) -- ([xshift=8mm]ring.north)
			node[midway, right, font=\scriptsize] {写入};
			% 回读 (偏移 -8mm)，直接向上
			\draw[route] ([xshift=-8mm]ring.north) -- ([xshift=-8mm]cpu.south)
			node[midway, left, font=\scriptsize] {回读};
			
			% 告警反馈
			\draw[route] (ids.south) |- (ring.east)
			node[pos=0.7, right, font=\scriptsize] {告警元数据};
			
			% 时间同步
			\draw[sync] (gptp.south) -- (cpu.north) node[midway, right, font=\scriptsize] {主时钟};
			\coordinate (sync_pt) at (0, 3.2);
			\draw[sync] (sync_pt) -| (can.north);
			\draw[sync] (sync_pt) -| (eth.north);
			\draw[sync] (sync_pt) -| (ids.north);
			
		\end{tikzpicture}
	}
	\caption{系统部署架构图}
\end{figure}
\end{document}